\documentclass[11pt]{article}
\usepackage[margin=1in]{geometry}
\usepackage{amsmath,amssymb,amsthm}
\usepackage[hidelinks]{hyperref}
\newtheorem{theorem}{Theorem}
\title{Cycle Packing at bounded degree and large girth}
\author{Atlas}
\date{18 November 2026}
\begin{document}
\maketitle

In \textsc{Cycle Packing}, we ask whether an explicit simple undirected
graph \(G\) has at least the given number of pairwise vertex-disjoint
cycles. Write \(n=|V(G)|\) and \(N\) for the complete binary input length,
including the integer threshold and any supplied decomposition.
For \(t\ge0\) let \(\phi(t)=\max\{1,t\log_2(\max\{t,2\})\}\).
A computable nonnegative integer function \(h\) is effectively \(o(\phi(t))\)
if, for each positive rational \(\eta\), a computable cutoff \(T(\eta)\)
guarantees \(h(t)\le \eta\phi(t)\) for all \(t\ge T(\eta)\).

\hypertarget{cycle-packing-degree-six}{}
\begin{theorem}\label{cycle-packing-degree-six}
Assume deterministic \(\mathrm{FPT}\ne\mathrm W[1]\). For every fixed
\(0<\gamma<1\), every fixed \(C\ge0\) and every computable effectively
\(o(\phi(t))\) function \(h\), no deterministic algorithm decides
\textsc{Cycle Packing} in \(O(2^{h(t)}N^C)\) time on connected bipartite
graphs of maximum degree six and girth at least \(n^{1-\gamma}\).
This holds for intrinsic treewidth and separately for the width of
a supplied path decomposition, even when the decomposition is optimal
and treewidth equals pathwidth.
\end{theorem}

The historical \(t\log t\) scale under ETH comes from Lokshtanov,
Marx and Saurabh~\cite[Theorem 6.3]{survey} and Cygan et al.~
\cite[Theorem 14.19]{book}. We adapt the permutation host of
Cygan et al.~\cite[Appendix C.2 of the full version]{ccfull}
to handle arbitrary sets of cells. A further replacement bounds the
degree while preserving the optimum up to a known offset. A stretched
grid supplies exact width without introducing short cycles.


\section{From Clique to permutation hitting}

We use the \(\mathrm{W[1]}\)-hardness of Clique parameterized by the requested
size \(k\)~\cite{df}. For a graph with \(n\ge 2\) vertices, use \(k\) blocks of
\(\lceil\log_2 n\rceil\) Boolean variables to encode the chosen vertices.
Forbid each unused code in each block, and for each pair of blocks forbid every
pair of valid codes representing equal or nonadjacent vertices. Each prohibition
is a clause on the affected blocks, falsified exactly by the forbidden pattern.
The formula is satisfiable exactly when the graph has a \(k\)-clique. It has
\(m=k\lceil\log_2 n\rceil\) variables, and its complete length and construction
time are both bounded by \(F(k)n^a\) for computable \(F\) and a fixed \(a\).
We decide trivial Clique inputs directly.

Take an explicit CNF over \(m\) listed variables. Decide instances with fewer
than \(16\) variables or an empty clause directly, and likewise decide a
formula with no clauses. Put
\[
 b=\lfloor (\log_2 m)/2\rfloor,\qquad
 q=2^b,\qquad s=\lceil m/b\rceil,\qquad d=\max\{s,q\}.
\]
Pad the variable list to \(sb\) and partition it into \(s\) groups of size
\(b\). In row \(i\) of a \(d\times d\) table, the first \(q\) columns encode
assignments to group \(i\), for each \(i\le s\). Require a row selection to
hit the set of these \(q\) cells for each real row. For each clause, require
it to hit the set of all real cells whose group assignment satisfies a
literal of the clause. Thus a mapping selecting one column per row hits all
sets exactly when its real rows encode a satisfying assignment. Dummy
rows have no domain constraint. The set family is explicit and has polynomial
size in the formula length and \(m\). We have \(d=O(m/\log m)\).

Use the deterministic constructive coloring family of Lokshtanov, Marx, and
Saurabh~\cite[Lemma 2.8 and its proof]{lms}. It enumerates maps
\(f\colon[d]\times[d]\to[d+1]\) in time
\(2^{O(d\log\log(\max\{d,4\}))}d^{O(1)}\). For any mapping selecting one cell
per row, some \(f\) gives the selected cells distinct colors and gives each
selected cell a color absent from the other cells of its row. Indeed, the
graph formed by a clique on the selected cells and the stars joining each
selection to its row is the cactus-grid graph of that lemma.

For each \(f\), retain the cells whose colors are unique within their rows, and
replace the column of every retained cell by its color. Apply this operation to
every set and additionally require a hit on the retained cells of each row.
Add one unconstrained row to make a square of size \(D=d+1\).
For any permutation hitting these new sets, the row-color choice in each
original row has a unique preimage, which gives an original selection hitting
all old sets. Conversely, for a selection hitting all old sets, the map
\(f\) guaranteed by the coloring family retains every chosen cell and assigns
them distinct colors. Assign the unused color to the new dummy row. This
gives a permutation hitting all new sets. The number of queries and the time
to enumerate the family are \(2^{o(m)}\) up to a fixed polynomial factor in
the explicit formula length, with a computable sublinear exponent bound
\(O(m\log\log m/\log m)\).


\section{From permutation hitting to Cycle Packing}

We allow a multigraph temporarily, with a pair of parallel edges
forming a two-cycle. For ports \(v_1,\ldots,v_d\), add private vertices
\(u_1,\ldots,u_{d-1}\) and double the edges from \(u_i\) to \(v_i,v_{i+1}\).
For any omitted port \(v_j\), match the private vertices to all other
ports using \(u_i v_i\) for \(i<j\) and \(u_i v_{i+1}\) otherwise.
Each match makes a two-cycle. We also use a second type of gadget,
with one private vertex joined by a doubled edge to every port.

When private vertices have no neighbors elsewhere and port sets are
disjoint, normalize a maximum packing separately at each gadget.
For the first type, remove all cycles meeting a private vertex.
Their number is at most the number of private vertices and the number
of currently free ports, since each removed cycle uses both.
The matching covers as many free ports as possible up to \(d-1\),
so replace the removed cycles by at least as many two-cycles.
If there are still \(a\ge2\) ports occupied by the other cycles,
delete at most \(a-1\) such cycles leaving at most one occupied port.
Indeed, if some cycle uses at least two ports delete all touching
cycles, and otherwise retain one. Replacing the two-cycles by a
full matching gains at least \(a-1\), since at most \(d-a\)
were present. For the second type, replace the private vertex's
cycle by a two-cycle on one of its ports, or add a two-cycle on a
free port if the private vertex was unused. If all ports were
used and the private vertex unused, remove one touching cycle
to free its port. Previously installed two-cycles survive because
the port sets are disjoint. After normalizing, the first type
has one port available to other cycles and the second type has
all but one port available.

Let the permutation have order \(K\) and let \(S_1,\ldots,S_M\)
be its nonempty target sets of cells. An empty target set or an
empty family can be decided directly. Use row-end vertices \(P_i,Q_i\)
and column-end vertices \(A_j,B_j\). For each pair \((i,j)\) add
paths \(P_i,p_{ij},A_j\) and \(Q_i,q_{ij},B_j\).
The rail for row \(i\) connects \(P_i\) to \(Q_i\) through pairs
\(x^L_{i,s},y^L_{i,s}\) in increasing order of stage \(s\).
The rail for column \(j\) connects \(A_j\) to \(B_j\) through
\(x^R_{j,s},y^R_{j,s}\) in the same way. For every occurrence
\(e=(i,j,s)\) in a target set add
\[
 x^L_{i,s},a_e,c_e,x^R_{j,s}\qquad\text{and}\qquad
 y^L_{i,s},b_e,d_e,y^R_{j,s}
\]
as paths, and add the triangle \(a_e,b_e,z_e\).
Use the first gadget type on each row's \(p\) ports and separately
on its \(q\) ports. At each stage use the first type separately
on the \(c\) and \(d\) ports, and the second type on the \(z\) ports.
Let \(g=2K(K-1)+\sum_s(2|S_s|-1)\) count all private vertices
of these gadgets. The threshold is \(g+K+\sum_s|S_s|\).

\paragraph{Completeness.}
For a permutation \(f\) hitting all sets choose one cell
\((i,f(i))\) at each stage. Use the selector paths at the ends
of row \(i\) and column \(f(i)\). Together with their rails,
the two crossing paths of each chosen cell split that pair of
rails into one more cycle per chosen occurrence. If no cell is
chosen on the pair, the two end selectors close one cycle.
Thus these rail cycles contribute \(K+M\). Pack the local
triangle of every unchosen occurrence, for \(\sum_s(|S_s|-1)\)
additional cycles. The port capacities permit all \(g\) gadget
two-cycles, reaching the threshold.

\paragraph{Soundness.}
Normalize and remove the \(g\) gadget two-cycles. Let \(Z\)
consist of all \(P_i,Q_i,a_e,b_e\). Removing \(Z\) from the remaining
host gives a forest. Every vertex of \(Z\) has at most one neighbor
in each component of this forest. Every surviving cycle therefore
uses at least two vertices of \(Z\). As the number of cycles is
at least \(|Z|/2\), each cycle uses exactly two and all of \(Z\)
is used.

The cycle containing \(P_i\) uses its row rail and exactly one
\(p\) port, by capacity one. At that port's column vertex \(A_j\)
it must enter the column rail. Otherwise the cycle would reach
another row-end vertex \(P_{i'}\) through a second port. The row
components incident to \(P_i,P_{i'}\) are distinct, so the cycle
could not return between them without using a third vertex of
\(Z\). Distinct \(P_i\) thus use distinct \(A_j\) and define a
permutation \(f\). At the start of the first stage the past side
pairs the row and column rail ends of each \((i,f(i))\) by a
path with exactly one previously encountered vertex of \(Z\).

Suppose this property holds on entering stage \(s\). Some \(z_e\)
is unused by the host cycles, by the second gadget's capacity.
The covered \(a_e\) therefore uses at least one entry route,
either directly from its row or through \(c_e\) from its column.
If it also uses \(b_e\), its cycle has three vertices of \(Z\),
including the one already on that entering path. Thus it uses
both entry routes. These entries must be paired by the
same past path, or their two previously encountered vertices
of \(Z\) would give a third with \(a_e\). Hence \(e=(i,f(i),s)\)
witnesses a hit. The use of \(c_e\) excludes every other \(c\)
port in this stage. For another occurrence \(e'\), a traversal
of \(a_{e'}\) from its row entry would have to continue through
\(b_{e'}\) or through \(z_{e'}\) and \(b_{e'}\), again giving three
vertices of \(Z\). Thus \(a_{e'},b_{e'},z_{e'}\) form their local
triangle. The covered \(b_e\) uses the \(d_e\) path to pair the
exits of the chosen row and column, with itself as the new
vertex of \(Z\). All other paired rails continue through the
stage. This proves the invariant and every target-set hit.


Keep the \(4K\) rail ends throughout a path decomposition.
At each end-selector system process one row at a time, keeping
its at most \(2K\) ports and private vertices. In the middle keep
the \(4K\) rail vertices of the current stage, and of the next
stage across a stage transition. Process occurrences within the
stage in selector order for both sparse selectors.
Keep the preceding \(c,d\) ports until the next occurrence
and their shared selector-private vertices are processed.
Keep the \(z\)-selector center throughout the stage.
Only a constant number of other vertices are needed at any
time. The supplied width is \(O(K)\), with polynomially many
bags in the explicit permutation input. To obtain a simple
graph, replace each multigraph edge by a length-three path.
Every cycle projects to a multigraph cycle, including parallel
two-cycles, and this preserves the packing optimum. Extend the
path decomposition using bags containing the old endpoints and
at most two new vertices while processing an edge.

\section{Degree, girth, and exact width}

The simple host is connected when all target sets are nonempty:
the end-selector paths link all row and column rails and each
other vertex connects to a rail or its port. We bound the degree
of a simple host using the following replacement. Order its
edges by assigned covering bags of the supplied path decomposition.
For a vertex of degree \(d>0\), make a backbone path of length
\(d-1\) in incident-edge order. At position \(i\) attach a
port \(p_i\) to the backbone vertex and a companion \(x_i\) to
\(p_i\). Every old edge joins its two ports. For \(1\le i\le d-2\)
make a center \(u_i\) adjacent to \(p_j,x_j\) for
\(j\in\{i,i+1,i+2\}\). All degrees are at most six. Ignore
isolated vertices. For \(d\ge2\), if at most two ports are unavailable and
the available companions are free, match the \(d-2\) centers
to the first \(d-2\) available ports in order. Their indices
\(j_i\) satisfy \(i\le j_i\le i+2\), so this gives \(d-2\)
disjoint triangles. A degree-one gadget has no center and cannot
be traversed by a cycle. Any old packing lifts along backbones using
at most two ports per gadget and all these triangles, with
offset \(\sum_v\max(0,d(v)-2)\).

For the reverse, normalize a maximum packing one gadget at a time.
For \(d\ge3\), remove the \(c\) cycles touching its centers and let
\(a\) count ports used by the kept cycles. Companions have degree
one when centers are ignored, so each occupied port uses its
backbone edge and its external edge. Each kept cycle enters and
leaves equally often, and \(a\) is even. If \(a\le2\), the \(d-2\)
triangles pay for removing the \(c\le d-2\) cycles.
If \(a\ge4\), each removed cycle using exactly one center in the
gadget had to use a port there, since without ports a center and
its companions form a star. Such cycles use distinct ports outside
the kept packing. If \(c_1\le d-a\) denotes their number,
then \(c_1+2(c-c_1)\le d-2\), giving \(c\le d-1-a/2\).
Delete at most \(a/2-1\) additional cycles to leave at most two
ports occupied. If fewer than \(a/2\) cycles touch the gadget
delete all of them. Otherwise they each use two ports and one
can be kept. In both cases the \(d-2\) triangles pay for
all removals. Their vertices are local and do not disturb
triangles in gadgets already processed. After normalization
discard the offset triangles. At most two ports per gadget
remain occupied, and any visited gadget is traversed once.
Contract those traversals to recover disjoint cycles of the host.

For the width, process old edges in the assigned order, keeping
the newest backbone vertex and the newest two port-companion
pairs per active old vertex. For each edge introduce the next
backbone, port and companion at its two endpoints and their
centers for the latest triples. Cover all new edges, remove the
centers, and discard the old backbone and excess pairs.
All events for an old vertex fit in its original bag interval.
At most \(5(w+1)+8\) vertices are needed for old supplied width
\(w\). The output and certificate have polynomial length.


Replace every edge by a path of even length \(L\ge4\).
The graph stays connected, becomes bipartite and has girth
at least \(L\). The packing optimum is unchanged and the
decomposition extends with at most two extra vertices per bag.
Independently of \(L\), choose the least even \(q\ge2\) strictly larger
than the pre-subdivision width plus two, leaving
room for one more vertex in a bag. Take a \(q\) by \(q\) grid
and stretch each vertical edge to length \(L+1\). This remains
bipartite and has girth at least \(L\), since any cycle uses a
vertical path. Its packing optimum is \(q^2/4\), attained by
tiling the original grid by squares. No packing can have more
cycles since every cycle in the original grid uses at least
four vertices and subdivisions preserve packing.

The grid forces width at least \(q\). In the unexpanded grid
take all crosses consisting of one row and one column of the
top-left \((q-1)\) by \((q-1)\) subgrid. Add the bottom row and
the rightmost column with its bottom vertex removed. These
sets are each connected and pairwise touch, and a set meeting
all of them needs at least \(q+1\) vertices. Indeed, with fewer
than \(q-1\) vertices in the subgrid some row and column are
missed. In a tree decomposition, the subtrees of bags meeting
a connected set are connected. Pairwise touching forces these
subtrees to intersect pairwise, so the Helly property for
subtrees gives a bag meeting all sets. This yields the lower
bound for the original grid, and it is a minor of the stretched
grid. Contractions cannot increase treewidth, by identifying
the endpoints in all bags of a decomposition and using the
bag containing their edge to ensure a connected lifetime.

The stretched grid also has a supplied path decomposition of
width \(q\). Sweep rows from left to right, maintaining the
latest vertex in each column. To advance a column, introduce
the next vertex on its vertical path before forgetting its
predecessor. The previous column's new row vertex still covers
the horizontal edge. This needs at most \(q\) representatives
and one successor. In the initial row start at the top-left corner.

Join the stretched host to that corner by a bridge from an
internal subdivision vertex of the host. The endpoints have
degree two before joining. Choose the two orientations of the
bipartitions so the bridge preserves bipartiteness. No cycle
uses the bridge. Retain the host endpoint from its first bag
until the end, add a bridge bag, and start the grid certificate
at its top-left corner. Both intrinsic widths now equal \(q\)
and the supplied certificate attains it. Add \(q^2/4\) to the
threshold as the grid contribution.

Let \(H\) be the combined order before the final long subdivisions.
The degree bounds give \(n\le cHL\) for an absolute integer \(c\).
For the fixed \(\gamma\), choose an integer \(a\ge(1-\gamma)/\gamma\)
and take the least even \(L\ge\max(4,(cH)^a)\).
Then \(L^\gamma\ge(cH)^{1-\gamma}\), giving \(L\ge n^{1-\gamma}\).
All output lengths and construction costs stay polynomial per
permutation instance for this fixed \(\gamma\).

\section{Running time}

For the compact \(m\)-variable formula from a Clique instance
of size \(k\) and order \(n_0\), each output has \(t=q=O(m/\log m)\).
Consequently \(\phi(t)\le Am\) with an absolute integer \(A\).
A hypothetical solver at effectively \(o(\phi(t))\) rate has
exponent at most \(m/(4k)\) beyond its computable cutoff at
\(1/(4Ak)\). Below this cutoff its computable finite prefix
contributes only a factor depending on \(k\).
The effective disjunction overhead is at most \(2^{m/(4k)}\)
beyond a computable cutoff in \(m\), with another computable
parameter factor for smaller \(m\). All other factors are
fixed polynomials in the explicit source size, as are the
complete output sizes. Since \(m=k\lceil\log_2 n_0\rceil\), the
result is a deterministic fixed-parameter algorithm for
Clique by requested size, a contradiction.

\paragraph{Acknowledgment.}
The proof was found using ChatGPT.

\begin{thebibliography}{9}
\bibitem{df}
R. G. Downey and M. R. Fellows. Fixed-parameter tractability and
completeness II: On completeness for \(\mathrm W[1]\).
\emph{Theoretical Computer Science}, 141(1--2):109--131, 1995.
\href{https://doi.org/10.1016/0304-3975(94)00097-3}{DOI}.
\bibitem{lms}
D. Lokshtanov, D. Marx, and S. Saurabh.
Slightly superexponential parameterized problems.
\emph{SIAM Journal on Computing}, 47(3):675--702, 2018.
\href{https://doi.org/10.1137/16M1104834}{DOI}.
\bibitem{survey}
D. Lokshtanov, D. Marx, and S. Saurabh.
Lower bounds based on the Exponential Time Hypothesis.
\emph{Bulletin of the EATCS}, 105, 2011.
\bibitem{book}
M. Cygan et al. \emph{Parameterized Algorithms}. Springer, 2015.
\href{https://doi.org/10.1007/978-3-319-21275-3}{DOI}.
\bibitem{ccfull}
M. Cygan et al.
Solving connectivity problems parameterized by treewidth in single
exponential time. \emph{ACM Transactions on Algorithms}, 18(2), 2022.
\href{https://doi.org/10.1145/3506707}{DOI}.
The inspected appendix is in
\href{https://arxiv.org/abs/1103.0534v1}{arXiv:1103.0534v1}.
\end{thebibliography}
\end{document}



